\section{Open-source strategy}

\subsection{The data commons as the core asset}

The most valuable long-term output of the project may not be the hardware design or even the software --- it may be the \textbf{shared open dataset}. A large, well-structured, openly licensed dataset of plant growth under controlled conditions, tagged with species, sensor readings, yield outcomes, disease events, and energy consumption, would be a genuinely novel resource. Nothing comparable currently exists for small-scale CEA at the level of granularity this system produces.

This dataset has value to researchers, plant breeders, agronomists, and other engineers independently of any single product. It should be treated as a first-class project output, not a by-product.

\subsection{Community governance}

The project should establish governance structures from the outset to prevent the data commons from becoming a de-facto proprietary asset:

\begin{itemize}
  \item The dataset is hosted on infrastructure that the community can mirror.
  \item Model weights are published with provenance (which dataset version, which training run).
  \item The species database is a separately maintained community project.
  \item Ingest schema changes require a deprecation period and community discussion.
\end{itemize}

\subsection{Self-hostability}

The entire cloud stack must be deployable by any community member with a modest VPS. The deployment target is \texttt{docker compose up} from a single repository. This prevents the project from creating a dependency on a single hosted service and ensures continuity if the founding contributors step back.

% -------------------------------------------------------

\section{Impact assessment}

The following is a summary evaluation of the project's potential impact across the eleven problems identified in Section~\ref{sec:problems}.

\begin{center}
\begin{tabular}{p{4.5cm}p{3cm}p{2cm}p{4cm}}
\toprule
\textbf{Problem} & \textbf{Impact type} & \textbf{Confidence} & \textbf{Key condition} \\
\midrule
Land use / habitat loss & Direct, strong & High & Deployment reaches meaningful scale \\
Biodiversity loss & Direct, strong & High & As above \\
Soil erosion & Bypass (soilless) & High & Not a remediation of damaged land \\
Fertiliser / eutrophication & Direct, strong & High & Closed-loop nutrient management implemented correctly \\
Pesticide pollution & Direct, strong & High & Enclosed environment maintained; fungal risks managed \\
Monoculture & Enables diversity & Medium & Depends on user and community crop choices \\
Freshwater depletion & Direct, strong & High & Hydroponic/aeroponic systems required \\
GHG emissions & Conditional & Medium & Conditional on low-carbon electricity \\
Transport / food waste & Direct, strong & High & Local deployment required \\
Supply-chain concentration & Direct, strong & High & Deployment breadth required \\
Dietary diversity & Direct & Medium & Species recommender + user education \\
\textbf{Energy intensity (new risk)} & \textbf{Risk introduced} & \textbf{High} & \textbf{Mitigated by recommender; not eliminated} \\
\bottomrule
\end{tabular}
\end{center}

% -------------------------------------------------------

\section{Future work}

\begin{itemize}
  \item Integration with home energy management systems (excess solar $\to$ priority grow-light scheduling).
  \item Federated learning implementation to allow model improvement without raw data leaving the node.
  \item Aquaponics extension: closed-loop fish + plant system for protein production.
  \item Seed library integration: linking the species database to open seed repositories.
  \item Economic model: cost-per-kg analysis across species, climate zones, and hardware configurations to help users understand the financial viability of their setup.
\end{itemize}
